\documentclass[12pt]{article}
\usepackage{../../ficha}

\begin{document}

\ficha{Frações Impróprias e Numerais Mistos}

\instrucoes{Resolve sem calculadora e apresenta as frações na forma irredutível.}

\secao{Frações Próprias e Impróprias}{Classifica cada fração como própria, imprópria ou aparente (igual a um número inteiro).}
\begin{exercicios}
  \item $\dfrac{3}{5}$
  \item $\dfrac{7}{4}$
  \item $\dfrac{12}{3}$
  \item $\dfrac{9}{10}$
  \item $\dfrac{15}{8}$
  \item $\dfrac{20}{5}$
  \item $\dfrac{11}{11}$
  \item $\dfrac{25}{6}$
  \item Escreve três frações impróprias com denominador 7.
  \item Escreve três frações aparentes com denominador 4.
\end{exercicios}

\secao{De Fração Imprópria para Numeral Misto}{Escreve cada fração imprópria na forma de numeral misto.}
\begin{exercicios}
  \item $\dfrac{7}{2}$
  \item $\dfrac{11}{3}$
  \item $\dfrac{9}{4}$
  \item $\dfrac{17}{5}$
  \item $\dfrac{23}{6}$
  \item $\dfrac{29}{8}$
  \item $\dfrac{40}{7}$
  \item $\dfrac{51}{10}$
  \item $\dfrac{35}{12}$
  \item $\dfrac{100}{9}$
\end{exercicios}

\secao{De Numeral Misto para Fração Imprópria}{Escreve cada numeral misto na forma de fração imprópria.}
\begin{exercicios}
  \item $1\,\dfrac{1}{2}$
  \item $2\,\dfrac{1}{3}$
  \item $3\,\dfrac{3}{4}$
  \item $4\,\dfrac{2}{5}$
  \item $5\,\dfrac{1}{6}$
  \item $2\,\dfrac{7}{8}$
  \item $6\,\dfrac{3}{10}$
  \item $7\,\dfrac{2}{3}$
  \item $10\,\dfrac{1}{4}$
  \item $3\,\dfrac{5}{12}$
\end{exercicios}

\secao{Representação e Comparação}{Responde a cada questão.}

\begin{center}
\begin{tikzpicture}[x=1.4cm]
  \draw[-{Latex}] (-0.2,0) -- (4.4,0);
  \foreach \i in {0,...,4} {\draw (\i,0.15) -- (\i,-0.15); \node[below] at (\i,-0.15) {$\i$};}
  \foreach \i in {0,...,15} \draw (\i/4,0.07) -- (\i/4,-0.07);
  \fill (1.25,0) circle (2pt) node[above=2pt] {$A$};
  \fill (2.75,0) circle (2pt) node[above=2pt] {$B$};
  \fill (3.5,0) circle (2pt) node[above=2pt] {$C$};
\end{tikzpicture}
\end{center}

\begin{exercicios}
  \item Indica, como fração imprópria e como numeral misto, os números representados por $A$, $B$ e $C$.
  \item Entre que dois números inteiros consecutivos está $\dfrac{19}{5}$?
  \item Compara usando $<$, $>$ ou $=$: $\;\dfrac{9}{4} \;\square\; 2\,\dfrac{1}{4}$
  \item Compara usando $<$, $>$ ou $=$: $\;\dfrac{13}{3} \;\square\; 4\,\dfrac{1}{2}$
  \item Compara usando $<$, $>$ ou $=$: $\;3\,\dfrac{2}{5} \;\square\; \dfrac{16}{5}$
  \item Ordena por ordem crescente: $\dfrac{7}{2},\ 3\,\dfrac{1}{4},\ \dfrac{10}{3},\ 3$.
  \item Escreve um numeral misto entre $\dfrac{5}{2}$ e $\dfrac{11}{4}$.
  \item Qual é o maior: $\dfrac{22}{7}$ ou $3\,\dfrac{1}{8}$?
\end{exercicios}

\secao{Operações com Numerais Mistos}{Calcula e apresenta o resultado como numeral misto (quando possível).}
\begin{exercicios}
  \item $1\,\dfrac{1}{2} + 2\,\dfrac{1}{4}$
  \item $3\,\dfrac{2}{3} + 1\,\dfrac{1}{2}$
  \item $5\,\dfrac{3}{4} - 2\,\dfrac{1}{4}$
  \item $4\,\dfrac{1}{3} - 1\,\dfrac{2}{3}$
  \item $2\,\dfrac{1}{2} \times 4$
  \item $1\,\dfrac{1}{3} \times 2\,\dfrac{1}{4}$
  \item $3\,\dfrac{3}{5} : 3$
  \item $2\,\dfrac{1}{4} : 1\,\dfrac{1}{2}$
  \item $6 - 2\,\dfrac{5}{6}$
  \item $\left(1\,\dfrac{1}{2}\right)^2$
\end{exercicios}

\secao{Desafios}{Resolve os problemas.}
\begin{exercicios}
  \item Uma receita leva $2\,\dfrac{1}{4}$ chávenas de farinha. Quantas chávenas são precisas para fazer a receita 3 vezes?
  \item Uma fita tem $7\,\dfrac{1}{2}$\,m. Quantos pedaços de $1\,\dfrac{1}{4}$\,m se podem cortar?
  \item O Pedro correu $3\,\dfrac{2}{5}$\,km de manhã e $2\,\dfrac{3}{4}$\,km à tarde. Quantos quilómetros correu no total?
  \item Um depósito de $12$ litros tem $8\,\dfrac{2}{3}$ litros de água. Quantos litros faltam para o encher?
  \item Uma pizza foi dividida em 8 fatias iguais. Comeram-se 19 fatias de várias pizzas iguais. Quantas pizzas se comeram, em numeral misto?
  \item Encontra o número que somado com $1\,\dfrac{5}{6}$ dá $4\,\dfrac{1}{3}$.
  \item Um terreno retangular mede $5\,\dfrac{1}{2}$\,m por $3\,\dfrac{1}{3}$\,m. Calcula a área.
  \item Qual é o menor numeral misto, com denominador 7, maior do que $\dfrac{30}{7}$?
\end{exercicios}

\end{document}
